\section{What the Config Calls a Subgraph}
\label{sec:ch17-what-the-config-calls-a-subgraph}

\index{router execution config!what it records about a subgraph|(}%
Section~\ref{sec:ch17-the-one-thing-it-asks-permission-for} ends on an
asymmetry that looks arbitrary until you go and read what the composer
produces. Then it stops being arbitrary, because the file it writes has no
place to put an owner and every place to put a route.

\index{router execution config!a subgraph is three properties}%
Chapter~\ref{ch:composition} pulled this file apart the first time. A subgraph
in it is one object, and this is the whole of it:

\begin{minted}[fontsize=\footnotesize,breakanywhere]{json}
{ "id": "2", "name": "ratings", "routingUrl": "http://localhost:5003/graphql" }
\end{minted}

\index{ownership!not a property of anything in the file}%
A name and an address. No team, no owner, no contact, no repository. Beside it
sits the datasource entry, which carries the entity keys that service can be
entered through, the fields it declares a \code{@requires} on, and two lists of
type names and field names. Every one of those is a statement about where a
value can be fetched from. \emph{Whose} it is is not a property of anything in
the file, which is why the two edits in
section~\ref{sec:ch17-the-field-nobody-approved} left it looking exactly as
healthy as it looked before.

\subsection{The Question the Composer Is Asking}

\index{composition!every directive is about a route}%
Go back through the five directives this book uses to describe a seam and the
pattern holds across all of them. \code{@key} says which fields let a service
be entered at a type. \code{@external} says a declaration is a borrowed copy
that this service cannot answer. \code{@requires} says a field cannot be
resolved until another service has produced something first. \code{@shareable}
says more than one service can answer, and \code{@override} says the one that
answers is now this one. Five directives, and all five are answers to
\emph{where does the router go for this value}. The book uses two others that
are not, \code{@authenticated} from chapter~\ref{ch:entities-authorization} and
\code{@inaccessible} in section~\ref{sec:ch17-writing-it-down-somewhere-else},
and both are about whether a client may have the value rather than where it
comes from. Neither is about whose it is either.

\index{ownership!routing and ownership overlap and are not the same}%
Which is not a criticism. That question has to be answered and no other tool
in the stack can answer it. The trouble is that routing and ownership overlap
enough to be confused for one another, and they come apart at exactly the
place this chapter is about. With the initially unmarked declarations in
this experiment, two services claiming the field stop composition until
both documents permit sharing. That is a schema requirement, not evidence
that both teams approved the change today. An already-shareable field needs
no fresh edit from its existing owner. An override is different again: the
new declaration selects one route without requiring the old document to
change. Neither mechanism records business approval.

\index{figure!two ways a field changes hands}%
Figure~\ref{fig:ch17-two-ways-to-move-a-field} is the two routes to the same
routing-table entry.

\begin{figure}[htbp]
  \centering
  % Two ways abstract moves from one subgraph to another, measured against wgc
% 0.129.9 and the book's four exported subgraph documents. Above, two teams
% share an initially unmarked field: both declarations need editing. A field
% already declared shareable would not require a fresh old-owner edit.
% Below, an override changes the route without editing the old declaration.
% Both preserve the client-facing field. Neither records business approval.
%
% Same idiom as figures/tikz/ch01-one-field.tex, which fixes it: x=1mm/y=1mm,
% ticks above a single-weight axis, event nodes anchored south with manual
% line breaks rather than relying on the text width to wrap. The axis, tick,
% event, span and spanlabel styles, and the outcome style's placement of both
% lanes' outcome nodes at one shared x position, are copied from
% figures/tikz/ch16-two-places-a-change-can-fail.tex, which in turn takes the
% outcome style from figures/tikz/ch15-the-order-decides.tex.
%
% Two departures from ch16, both because the two lanes here are not the same
% shape as ch16's:
%
%   1. Lane A's cross does not end the lane. In ch16 a cross marks a walk
%      reaching a point it cannot pass, and the axis stops there because
%      nothing after it happened. Here the composer's refusal is not the end
%      of the story, it is the mechanism doing its job: it sends the change
%      back, both teams mark the field shareable, and the second compose
%      succeeds. So the axis is drawn as one continuous stroke from the first
%      event to the last, and the cross is laid across it exactly the way an
%      ordinary tick already crosses the axis without breaking it. Only the
%      shape of the mark changes, not the line beneath it.
%
%   2. Lane B carries no cross mark at all. Every other two-lane timeline in
%      this book puts a cross on both lanes; this figure puts one only on
%      lane A, because nothing on lane B goes wrong anywhere along it, and
%      that absence is the second half of the finding, not an oversight.
%
% Lane B's axis also stops noticeably short of the shared outcome column,
% the same device ch16 uses for its refused lane: the blank run-in says
% nobody had anything further to do.
%
% Spacing note. The first draft put five events on lane A at 22mm apart and
% the labels collided on the page, which the printed proof caught and the
% build log did not, since a node overlapping its neighbour is not an
% overfull box. Two things fixed it: the two shareable edits became one
% event, because what the lane argues is that both documents have to move
% rather than in which order, and the remaining ticks went to 32mm with the
% event text width at 30mm. Anything wider than the tick spacing collides,
% so a later edit that lengthens a label has to move the ticks with it.
%
% No quotation marks anywhere in this file, including this comment block, per
% the note at the head of ch07-the-walk.tex and repeated in
% ch15-the-order-decides.tex.
%
% Bare tikzpicture (house style): the figure environment, caption and label
% live at the call site in the section file.
\begin{tikzpicture}[
  x=1mm, y=1mm,
  every node/.append style={font=\sffamily\scriptsize},
  axis/.style={draw, line width=0.4pt,
    -{Straight Barb[length=1.4mm, width=1.6mm]}},
  tick/.style={draw, line width=0.4pt},
  event/.style={anchor=south, align=center, text width=30mm, inner sep=1pt},
  outcome/.style={anchor=west, align=left, text width=32mm, inner sep=1pt},
  span/.style={draw, line width=0.4pt},
  spanlabel/.style={anchor=north, align=center, inner sep=2pt},
  lane/.style={anchor=west, font=\sffamily\scriptsize\itshape},
]

% ------------------------------------------------------------------- lane A
\node[lane] at (0,20) {A. sharing an initially unmarked field};

% One continuous stroke from the first event to the last: the refusal at
% x=46 is not where the walk stops, so the axis does not stop there either.
\draw[axis] (0,0) -- (112,0);

\draw[tick] (14,-1.5) -- (14,1.5);
\draw[tick] (44.8,-1.5) -- (47.2,1.5);
\draw[tick] (44.8,1.5) -- (47.2,-1.5);
\draw[tick] (78,-1.5) -- (78,1.5);
\draw[tick] (106,-1.5) -- (106,1.5);

\node[event, text width=28mm] at (14,3)
  {ratings adds abstract\\to its own document};
\node[event, text width=26mm] at (46,3)
  {wgc router compose\\exit 1, refused};
\node[event, text width=30mm] at (78,3)
  {shareable is added\\to both documents};
\node[event, text width=24mm] at (106,3)
  {wgc router\\compose, exit 0};

\draw[span] (30,-4) -- (30,-7) -- (94,-7) -- (94,-4);
\node[spanlabel, text width=44mm] at (62,-8)
  {both declarations must permit sharing; an existing shareable declaration
   would need no fresh edit};

\node[outcome] at (120,0) {both documents moved;\\approval is not recorded.};

% ------------------------------------------------------------------- lane B
\begin{scope}[shift={(0,-50)}]
  \node[lane] at (0,20) {B. overriding needs no old-owner edit};

  % Stops well short of the outcome column: nothing happens after the field
  % is served, and nobody is asked along the way, which is the point of the
  % lane, so the blank run-in is doing some of the arguing.
  \draw[axis] (0,0) -- (98,0);
  \foreach \x in {14,46,82}{\draw[tick] (\x,-1.5) -- (\x,1.5);}

  \node[event, text width=28mm] at (14,3)
    {ratings adds abstract\\with @override\\(from: sessions)};
  \node[event, text width=30mm] at (46,3)
    {wgc router compose\\exit 0, nothing printed};
  \node[event, text width=26mm] at (82,3)
    {the router loads it\\and serves the field};

  \draw[span] (2,-4) -- (2,-7) -- (58,-7) -- (58,-4);
  \node[spanlabel, text width=44mm] at (30,-8)
    {override selects the new route without changing the old declaration};

  \node[outcome] at (120,0)
    {one document moved;\\approval is not recorded.};
\end{scope}

\end{tikzpicture}

  \caption{The same field ending up answerable by the ratings service, twice.
  Above, the initially unmarked field needs \code{@shareable} in both
  documents. Below, \code{@override} changes the route without an edit to
  the old owner's document. These are the edits this experiment required,
  not an approval workflow: an existing shareable declaration would remove
  the first lane's need for a fresh edit by that owner.}
  \label{fig:ch17-two-ways-to-move-a-field}
\end{figure}

\subsection{Where the Specifications Stand}

\index{ownership!absent from the specification too}%
This is not a gap in one composer. Apollo's directive reference lists no
directive that records a team, an owner or a contact. What it does list is
\code{@composeDirective}, whose whole job is to tell composition that
\enquote{all uses of a particular custom type system directive in the subgraph
schema should be preserved in the supergraph schema} rather than dropped on
the way through~\autocite{apollodirectives}. So the answer exists and it is
worth knowing, and it is also an exact statement of the position: if you want
an owner in your schema, you define the directive and the meaning yourself,
and the toolchain agrees to carry it without ever reading it.

\index{Composite Schemas specification!does not use the word}%
The specification chapter~\ref{ch:federation-model} declined to build on is
in the same place. Search the Composite Schemas draft for \emph{owner} or
\emph{ownership} and neither word occurs
anywhere~\autocite{compositeschemasspec}. A
specification that added an owner would have to say what an implementation
must do differently when it reads one, and there is no answer to that: an
owner is not a routing fact, so a router has nothing to do with it.

\index{ownership!lives outside the toolchain}%
Which leaves the record where it started, in the heads of the people who work
on the graph, or in whatever they wrote down themselves.
Section~\ref{sec:ch17-writing-it-down-somewhere-else} is about where to put it.
First, though, the directive that does move a field, used properly, because
there is a right way to do what section~\ref{sec:ch17-the-field-nobody-approved}
did wrong.
\index{router execution config!what it records about a subgraph|)}%
